% Part: turing-machines 
% Chapter: undecidability 
% Section: verification

\documentclass[../../../include/open-logic-section]{subfiles}

\begin{document}

\olfileid[tr]{tur}{und}{ver} 
\olsection{Gösterimin Doğrulanması}

\begin{explain}
Gösterimimizin işlediğini doğrulamak için iki şey kanıtlamalıyız. Önce $M$,
$w$ girdisinde duruyorsa $!T(M,w)\lif !E(M,w)$ tümcesinin geçerli olduğunu
göstermeliyiz. Ardından tersini, yani $!T(M,w)\lif !E(M,w)$ geçerliyse
$M$'nin $w$ girdisinde çalıştırıldığında gerçekten sonunda durduğunu
göstermeliyiz.

Bunları kanıtlama stratejileri birbirinden çok farklıdır. İlk sonuç için
birinci dereceden mantığın !!a{sentence}, yani $!T(M,w)\lif !E(M,w)$'nin
geçerli olduğunu göstermeliyiz. Bunu yapmanın en kolay yolu !!a{derivation}
vermektir. Ne var ki kanıtımızın bütün $M$ ve $w$ için işlemesi gerekir;
dolayısıyla tek bir !!{sentence} yoktur ki onun için !!a{derivation} verelim;
aksine sonsuz sayıda tümce vardır. Bu nedenle yapabileceğimiz en iyi şey,
$M$ ile
$w$ nasıl olursa olsun ve $M$'nin $w$ girdisinde durması kaç adım sürerse
sürsün, $!T(M,w)\lif !E(M,w)$ için !!a{derivation} bulunacağını tümevarımla
kanıtlamaktır.

Doğal olarak tümevarımımız, $M$'nin bir durma konfigürasyonuna ulaşmadan önce
attığı adımların sayısı üzerine ilerleyecektir. Tümevarımlı kanıtımızda
$M$'nin $w$ girdisindeki çalışmasının her $n$ adımı için
$!T(M,w)\Entails !C(M,w,n)$ olduğunu ortaya koyacağız; burada $!C(M,w,n)$,
$M$'nin $w$ üzerinde $n$ adım çalıştırıldıktan sonraki konfigürasyonunu doğru
biçimde betimler. $M$, $w$ girdisinde örneğin $n$ adım sonra duruyorsa
$!C(M,w,n)$ bir durma konfigürasyonunu betimler. $!C(M,w,n)$ bir durma
konfigürasyonunu betimlediğinde $!C(M,w,n)\Entails !E(M,w)$ olduğunu da
göstereceğiz. Dolayısıyla $M$, $w$ girdisinde durursa bir $n$ için $n$ adım
sonra bir durma konfigürasyonunda olacaktır. Böylece $!C(M,w,n)$ bir durma
konfigürasyonunu betimlerken $!T(M,w)\Entails !C(M,w,n)$ ve bu durumda
$!C(M,w,n)\Entails !E(M,w)$ olduğundan $!T(M,w)\Entails !E(M,w)$, yani
$\Entails !T(M,w)\lif !E(M,w)$ elde ederiz.

Tersini kanıtlama stratejisi çok farklıdır. Burada
$\Entails !T(M,w)\lif !E(M,w)$ varsayımından $M$'nin $w$ girdisinde durduğunu
kanıtlamalıyız. Varsayımdan $!T(M,w)\Entails !E(M,w)$, yani $!T(M,w)$'nin
doğru olduğu her !!{structure} içinde $!E(M,w)$'nin doğru olduğu çıkar.
Bu yüzden $!T(M,w)$'nin doğru olduğu bir !!a{structure}~$\Struct{M}$ betimleriz:
Evreni $\Nat$ olacak ve bütün $\Obj Q_q$ ile $\Obj S_\sigma$ yorumları,
$M$'nin $w$ girdisindeki çalışmasının konfigürasyonlarıyla verilecektir.
Örneğin $\Sat{M}{\Obj Q_q(\num{m},\num{n})}$ ancak ve ancak $M$, $w$
girdisinde $n$ adım çalıştırıldıktan sonra $m$. kareyi tararken $q$
durumundaysa geçerli olacaktır. Varsayım gereği
$!T(M,w)\Entails !E(M,w)$ ve kuruluş gereği $\Sat{M}{!T(M,w)}$ olduğundan
$\Sat{M}{!E(M,w)}$ olur. Ancak $\Sat{M}{!E(M,w)}$ ancak ve ancak öyle bir
$n\in\Domain{M}=\Nat$ varsa ki $M$, $w$ girdisinde
çalıştırıldığında $n$ adım sonra bir durma konfigürasyonundaysa geçerlidir.
\end{explain}

\begin{defn} 
$!C(M,w,n)$ şu !!{sentence} olsun:
\[ 
\Obj Q_q(\num{m}, \num{n}) \land \Obj S_{\sigma_0}(\num{0}, \num{n})
\land \dots \land \Obj S_{\sigma_k}(\num{k}, \num{n}) \land
\lforall[x][(\num{k} < x \lif \Obj S_\TMblank(x, \num{n}))]
\] 
Burada $q$, $n$ anında $M$'nin durumu; $m$, $n$ anında $M$'nin taradığı
kare; $\sigma_i$, $0\le i\le k$ için $n$ anında $i$. karenin içerdiği simge
ve $k$ de $0$ anında şeridin en sağdaki boş olmayan karesiyle $n$ adım içinde
şerit kafasının ziyaret ettiği en sağdaki kareden daha büyük olanıdır.
\end{defn}

\begin{lem}\ollabel{lem:halt-config-implies-halt}
$M$, $w$ girdisinde çalıştırıldığında $n$ adım sonra bir durma
konfigürasyonundaysa $!C(M,w,n)\Entails !E(M,w)$.
\end{lem}

\begin{proof}
$M$'nin $w$ girdisinde $n$ adım sonra durduğunu varsayalım. Öyle bir $q$
durumu, $m$ karesi ve $\sigma$ simgesi vardır ki:
\begin{enumerate} 
\item $n$ adım sonra $M$, üzerinde $\sigma$ bulunan $m$. kareyi tararken
  $q$ durumundadır.
\item Geçiş fonksiyonu $\delta(q,\sigma)$ tanımsızdır.
\end{enumerate}
$!C(M,w,n)$ bu konfigürasyonun betimidir ve
$\Obj Q_q(\num{m},\num{n})$ ile
$\Obj S_\sigma(\num{m},\num{n})$ bağlaçlarını içerir. Bu bağlaçlar birlikte
$!E(M,w)$'yi gerektirir:
\[
\lexists[x][\lexists[y][(\bigvee_{\tuple{q, \sigma} \in
      X}(\Obj Q_q(x, y) \land \Obj S_\sigma(x, y)))]]
\]
çünkü $\tuple{q,\sigma}\in X$ olduğundan
$\Obj Q_q(\num{m},\num{n})\land\Obj S_\sigma(\num{m},\num{n})
\Entails \bigvee_{\tuple{q',\sigma'}\in X}
(\Obj Q_{q'}(\num{m},\num{n})\land
\Obj S_{\sigma'}(\num{m},\num{n}))$.
\end{proof}

\begin{explain}
Dolayısıyla $M$, $w$ girdisinde durursa öyle bir $n$ vardır ki
$!C(M,w,n)\Entails !E(M,w)$. Şimdi her $n$ anı için
$!T(M,w)\Entails !C(M,w,n)$ olduğunu göstereceğiz.
\end{explain}

\begin{lem}
\ollabel{lem:config}
Her $n$ için, $M$ $n$ adım sonunda henüz durmamışsa
$!T(M,w)\Entails !C(M,w,n)$.
\end{lem}

\begin{proof}
Tümevarım tabanı: $n=0$ ise $!C(M,w,0)$'ın bağlaçları aynı zamanda
$!T(M,w)$'nin bağlaçlarıdır; dolayısıyla onun tarafından gerektirilirler.

Tümevarım varsayımı: $M$, $n$. adımdan önce durmamışsa
$!T(M,w)\Entails !C(M,w,n)$. $!C(M,w,n)$ bir durma konfigürasyonunu
betimlemediği sürece $!T(M,w)\Entails !C(M,w,n+1)$ olduğunu göstermeliyiz.

$n\ge 0$ olduğunu ve $w$ üzerinde başlatılan $M$'nin $n$ adım sonra $m$.
kareyi tararken $q$ durumunda olduğunu varsayalım. $M$, $n$ adım sonra
durmadığına göre $M$'nin programında şu üç biçimden birinde bir yönerge
bulunmalıdır:

\begin{enumerate} 
\item \ollabel{right} $\delta(q, \sigma) = \tuple{q', \sigma', \TMright}$

\item \ollabel{left} $\delta(q, \sigma) = \tuple{q', \sigma', \TMleft}$

\item \ollabel{stay} $\delta(q, \sigma) = \tuple{q', \sigma', \TMstay}$
\end{enumerate}

Bu üç durumu sırayla ele alacağız.

\begin{enumerate} 
\item \olref{right} biçiminde bir yönerge bulunduğunu varsayalım.
  \olref[rep]{defn:tm-descr}\olref[rep]{rep-right} gereği bu,
\begin{align*} 
& \lforall[x][\lforall[y][((\Obj Q_{q}(x,
  y) \land \Obj S_\sigma(x, y)) \lif {}]]\\
& \qquad (\Obj Q_{q'}(x',
  y') \land \Obj S_{\sigma'}(x, y') \land !A(x, y)))  
\intertext{tümcesinin $!T(M,w)$'nin bir bağlacı olduğu anlamına gelir. Bu,
  şu !!{sentence} gerektirir (tümel örnekleme; $x$ yerine $\num{m}$ ve
  $y$ yerine $\num{n}$):}
& (\Obj Q_{q}(\num{m}, \num{n}) \land \Obj S_{\sigma}(\num{m},
\num{n})) \lif {}\\
& \qquad (\Obj Q_{q'}(\num{m}', \num{n}') \land
\Obj S_{\sigma'}(\num{m}, \num{n}') \land !A(\num{m}, \num{n})).
\intertext{Tümevarım varsayımına göre $!T(M,w)\Entails !C(M,w,n)$, yani}
& \Obj Q_q(\num{m}, \num{n}) \land \Obj S_{\sigma_0}(\num{0}, \num{n})
\land \dots \land \Obj S_{\sigma_k}(\num{k}, \num{n}) \land \\
& \qquad\lforall[x][(\num{k} < x \lif \Obj S_\TMblank(x, \num{n}))]\\
\intertext{$n$ adım sonra $m$. kare $\sigma$ içerdiği için karşılık gelen
  bağlaç $\Obj S_\sigma(\num{m},\num{n})$'dir; dolayısıyla bu şunu gerektirir:}
& \Obj Q_{q}(\num{m}, \num{n}) \land \Obj S_{\sigma}(\num{m},
   \num{n})
\intertext{Şimdi şunu elde ederiz:}
& \Obj Q_{q'}(\num{m}', \num{n}') \land \Obj S_{\sigma'}(\num{m},
  \num{n}') \land {}\\
& \qquad\bigwedge_{\substack{0\le i\le k\\i\ne m}}
  \Obj S_{\sigma_i}(\num{i},\num{n}') \land {}\\
& \qquad \lforall[x][(\num{k} < x \lif \Obj S_\TMblank(x, \num{n}'))]
\end{align*}
Şöyle: İlk satır önceki koşulun art bileşeninden doğrudan modus ponens ile
çıkar. Orta satırdaki her bağlaç, $\Obj
S_{\sigma_m}(\num{m},\num{n}')$ dışında, $!C(M,w,n)$ içindeki karşılık gelen
bağlaçla $!A(\num{m},\num{n})$'den çıkar.

$m<k$ ise $!T(M,w)\Proves \num{m}<\num{k}$
(\olref[rep]{prop:mlessk}) ve $<$ bağıntısının geçişliliğinden
$\lforall[x][(\num{k}<x\lif\num{m}<x)]$ elde ederiz. $m=k$ ise
$\lforall[x][(\num{k}<x\lif\num{m}<x)]$ yalnızca mantıktan çıkar. Son satır
da $!C(M,w,n)$ içindeki karşılık gelen bağlaç,
$\lforall[x][(\num{k}<x\lif\num{m}<x)]$ ve $!A(\num{m},\num{n})$'den çıkar.
$m<k$ ise bu zaten $!C(M,w,n+1)$'dir.

Şimdi $m=k$ olduğunu varsayalım. Bu durumda $n+1$ adım sonra şerit kafası,
artık ziyaret edilmiş en sağdaki kare olan $k+1$. kareyi de ziyaret etmiştir.
Dolayısıyla $!C(M,w,n+1)$ yeni bir
$\Obj S_\TMblank(\num{k}',\num{n}')$ bağlacı içerir ve son bağlaç
$\lforall[x][(\num{k}'<x\lif\Obj S_\TMblank(x,\num{n}'))]$ olur. Bu iki
!!{sentence}s de gerektirildiğini doğrulamalıyız.

$\lforall[x][(\num{k}<x\lif\Obj S_\TMblank(x,\num{n}'))]$ zaten elimizdedir.
Özellikle bundan $\num{k}<\num{k}'\lif
\Obj S_\TMblank(\num{k}',\num{n}')$ elde edilir.
$\lforall[x][x<x']$ aksiyomundan $\num{k}<\num{k}'$ elde ederiz. Modus
ponens ile $\Obj S_\TMblank(\num{k}',\num{n}')$ çıkar.

Ayrıca $!T(M,w)\Proves \num{k}<\num{k}'$ olduğundan $<$ bağıntısının
geçişlilik aksiyomu bize
$\lforall[x][(\num{k}'<x\lif\Obj S_\TMblank(x,\num{n}'))]$ verir. (Bunun
doğrulanmasını alıştırma olarak bırakıyoruz.)

\item \olref{left} biçiminde bir yönerge bulunduğunu varsayalım.
  \olref[rep]{defn:tm-descr}\olref[rep]{rep-left} gereği
\begin{align*} 
& \lforall[x][\lforall[y][((\Obj Q_{q}(x', y) \land \Obj
    S_{\sigma}(x', y)) \lif {}]]\\
& \qquad (\Obj Q_{q'}(x, y') \land \Obj
  S_{\sigma'}(x', y') \land !A(x', y))) \land {}\\
& \lforall[y][((\Obj Q_{q}(\num{0}, y) \land \Obj S_{\sigma}(\num{0},
    y)) \lif {}]\\
& \qquad (\Obj Q_{q'}(\num{0}, y') \land \Obj S_{\sigma'}(\num{0}, y')
  \land !A(\num{0}, y)))
\intertext{!!{sentence}, $!T(M,w)$'nin bir bağlacıdır. $m>0$ ise $l=m-1$ (yani
  $m=l+1$) olsun. Yukarıdaki tümcenin ilk bağlacı şunu gerektirir:}
& (\Obj Q_{q}(\num{l}', \num{n}) \land \Obj S_{\sigma}(\num{l}', \num{n}))
\lif {} \\
& \qquad (\Obj Q_{q'}(\num{l}, \num{n}') \land \Obj S_{\sigma'}(\num{l}',
\num{n}') \land !A(\num{l}', \num{n}))
\intertext{Aksi durumda $l=m=0$ olsun ve ikinci bağlacın gerektirdiği şu
  !!{sentence} ele alınsın:}
& ((\Obj Q_{q}(\num{0}, \num{n}) \land \Obj S_{\sigma}(\num{0}, \num{n})) \lif {}\\
& \qquad   (\Obj Q_{q'}(\num{0}, \num{n}') \land \Obj S_{\sigma'}(\num{0}, \num{n}') \land
!A(\num{0}, \num{n}))) 
\intertext{Her iki tümce de şunu gerektirir:}
&  \Obj Q_{q'}(\num{l}, \num{n}') \land \Obj S_{\sigma'}(\num{m},
  \num{n}') \land {}\\
&\qquad \bigwedge_{\substack{0\le i\le k\\i\ne m}}
  \Obj S_{\sigma_i}(\num{i},\num{n}') \land {} \\
& \qquad  \lforall[x][(\num{k} < x
  \lif \Obj S_\TMblank(x, \num{n}'))]
\end{align*}
önceki durumdaki gibi çıkar. (İlk durumda
$\num{l}'\ident\num{l+1}\ident\num{m}$, ikinci durumda ise
$\num{l}\ident\num{0}$ olduğuna dikkat edin.) Bu tam olarak
$!C(M,w,n+1)$'dir.

\item \olref{stay} durumu alıştırma olarak bırakılmıştır.
\end{enumerate}
Her $n$ için $!T(M,w)\Entails !C(M,w,n)$ olduğunu gösterdik.
\end{proof}

\begin{prob}
\olref[tur][und][ver]{stay} durumunu
\olref[tur][und][ver]{lem:config} kanıtında tamamlayın.
\end{prob}

\begin{prob}
$\Obj S_{\sigma_i}(\num{i},\num{n})$ ile $!A(m,n)$'den
$\Obj S_{\sigma_i}(\num{i},\num{n}')$ için !!a{derivation} verin
($i\neq m$, yani $i<m$ ya da $m<i$ olduğunu varsayın).
\end{prob}

\begin{prob}
$\lforall[x][(\num{k}<x\lif\Obj S_\TMblank(x,\num{n}'))]$,
$\lforall[x][x<x']$ ve
$\lforall[x][\lforall[y][\lforall[z][ ((x < y \land y < z) \lif x <
      z)]]]$ tümcelerinden
$\lforall[x][(\num{k}'<x\lif\Obj S_\TMblank(x,\num{n}'))]$ için
!!a{derivation} verin.)
\end{prob}


\begin{lem}
\ollabel{lem:valid-if-halt}
$M$, $w$ girdisinde durursa $!T(M,w)\lif !E(M,w)$ geçerlidir.
\end{lem}

\begin{proof}
\olref{lem:config} gereği her $n$ anında $M$'nin konfigürasyonunun
$!C(M,w,n)$ betiminin $!T(M,w)$ tarafından gerektirildiğini biliyoruz.
$M$'nin $k$ adım sonra durduğunu varsayalım. O anda bir $m\in\Nat$ için $m$.
kareyi tarıyor olacaktır. O zaman $!C(M,w,k)$, $M$'nin bir durma
konfigürasyonunu betimler; yani $\delta(q,\sigma)$ tanımsızken hem
$\Obj Q_q(\num{m},\num{k})$ hem de
$\Obj S_\sigma(\num{m},\num{k})$ ifadelerini bağlaç olarak içerir.
Dolayısıyla \olref{lem:halt-config-implies-halt} gereği
$!C(M,w,k)\Entails !E(M,w)$. Ancak
$!T(M,w)\Entails !C(M,w,k)$ olduğundan $!T(M,w)\Entails !E(M,w)$ ve böylece
$!T(M,w)\lif !E(M,w)$ geçerlidir.
\end{proof}


\begin{explain} 
İddiamızın doğrulanmasını tamamlamak için ters yönü de ortaya koymalıyız:
$!T(M,w)\lif !E(M,w)$ geçerliyse $M$, $w$ girdisinde başlatıldığında
gerçekten durur.
\end{explain}

\begin{lem}
\ollabel{lem:halt-if-valid}
$\Entails !T(M,w)\lif !E(M,w)$ ise $M$, $w$ girdisinde durur.
\end{lem}

\begin{proof}
Evreni $\Nat$ olan; $\Obj 0$'ı $0$, $\prime$ simgesini ardıl fonksiyonu,
$<$ simgesini küçüktür bağıntısı olarak ve $\Obj Q_q$ ile $\Obj S_\sigma$
yüklemlerini şöyle yorumlayan $\Lang L_M$-!!{structure}~$\Struct M$'yi ele
alın:
\begin{align*}
  \Assign{\Obj Q_q}{M} & =
\Setabs{\tuple{m, n}}{\begin{array}{ll}\text{$w$ üzerinde başlatılan $M$, $n$ adım sonra}\ \text{$m$. kareyi tararken $q$ durumundadır}
  \end{array}} \\
\Assign{\Obj S_\sigma}{M} & = \Setabs{\tuple{m, n}}{\begin{array}{ll}
\text{$w$ üzerinde başlatılan $M$'nin $n$ adım sonra}\ \text{$m$. karesi $\sigma$ simgesini içerir}
  \end{array}}
\end{align*}
Başka bir deyişle !!{structure}~$\Struct{M}$'yi, $w$ girdisinde başlatılan
$M$'nin gerçekten yaptıklarını adım adım betimleyecek biçimde kurarız.
Açıkça $\Sat{M}{!T(M,w)}$. $\Entails !T(M,w)\lif !E(M,w)$ ise
$\Sat{M}{!E(M,w)}$ de olur; yani
\[
\Sat{M}{\lexists[x][\lexists[y][(\bigvee_{\tuple{q, \sigma} \in
      X}(\Obj Q_q(x, y) \land \Obj S_\sigma(x, y)))]]}.
\]
$\Domain{M}=\Nat$ olduğundan $\delta(q,\sigma)$ tanımsız olan bazı $q$ ve
$\sigma$ için öyle bir $m$ ve $n\in\Nat$ bulunmalıdır ki
$\Sat{M}{\Obj Q_{q}(\num{m},\num{n})\land
\Obj S_{\sigma}(\num{m},\num{n})}$. $\Struct M$ tanımına göre bu, $w$
girdisinde başlatılan $M$'nin $n$ adım sonra $\sigma$ simgesini okurken $q$
durumunda ve geçiş fonksiyonunun tanımsız olduğu, yani $M$'nin durmuş olduğu
anlamına gelir.
\end{proof}

\end{document}
